\label{ch:chapter_finale}%
\chapter{Conclusioni e sviluppi futuri}

\section{Conclusioni}

La metodologia sviluppata ha consentito di ottenere una configurazione turbogetto--ramjet
compatibile con i requisiti assegnati nei due punti di progetto, integrando nello stesso
processo il dimensionamento termodinamico e geometrico dei principali componenti. Il
turbogetto, dimensionato per una spinta di 25 kN a $M=0.85$ e $h=12$ km, presenta un
TSFC di 0.135 kg/(N$\cdot$h) e un ingombro radiale massimo di 539 mm, con tutti i vincoli
aerodinamici e geometrici imposti in ottimizzazione rispettati. Il ramjet, dimensionato per 65 kN a $M=3.5$ e $h=25$ km,
presenta un TSFC di 0.231 kg/(N$\cdot$h), ampiamente entro il limite di missione, con un
diametro di efflusso di 1.57 m compatibile con il vincolo geometrico imposto.

%---------------------------------------
\section{Sviluppi futuri}

La struttura modulare della procedura di calcolo consente di aumentare
progressivamente il livello di dettaglio del modello, introducendo nuovi
sottosistemi oppure sostituendo quelli esistenti con formulazioni più
accurate. In seguito sono riportati possibili miglioramenti.

Un primo affinamento riguarda il modello termodinamico. Le proprietà dei fluidi sono state
valutate assumendo $c_p$ e $\gamma$ costanti e la combustione senza determinazione esplicita della
composizione dei prodotti. L'introduzione di proprietà dipendenti dalla temperatura e di un
modello di equilibrio chimico consentirebbe di aggiornare localmente entalpia, temperatura e
rendimento del ciclo.

La modellazione del flusso interno può essere estesa introducendo perdite viscose distribuite
nei condotti e perdite localizzate nelle transizioni geometriche. Un modello quasi-
monodimensionale, eventualmente affiancato da analisi CFD nei componenti più critici,
permetterebbe di valutare perdite di pressione totale, sviluppo dello strato limite, separazioni e
disuniformità del flusso nelle sezioni di interfaccia.

Per il compressore, il dimensionamento può essere esteso ai singoli stadi, aggiornando
localmente triangoli di velocità, grado di reazione, solidità e coefficienti di diffusione.
L'introduzione di correlazioni per perdite di profilo, flussi secondari e giochi di estremità
permetterebbe inoltre di verificare i margini rispetto a stallo e choking lungo l'intera macchina.

Per la turbina, il modello di raffreddamento può essere sostituito da una ripartizione stadio
per stadio della portata spillata tra statore e rotore, includendo raffreddamento interno,
film cooling e miscelazione con il flusso principale. Il conseguente bilancio termico locale
consentirebbe di stimare la temperatura del metallo e di introdurre direttamente
nell'ottimizzazione vincoli termici e aerodinamici sui singoli stadi.